\subsection{Zielsetzung}
\label{subsec:zielsetzung}

Ziel der Fallstudie ist eine strukturierte Untersuchung des Projekts
\enquote{Template-Projekte}. Betrachtet wird das darin entwickelte
Technologie-Template. Hierzu werden seine Architektur, die Aufteilung in
gemeinsam genutzte und profilspezifische Bestandteile sowie die unterstützenden
Entwicklungs- und Qualitätssicherungsprozesse analysiert. Die Betrachtung
richtet sich damit nicht nur auf die eingesetzten Technologien, sondern
insbesondere auf deren Zusammenwirken innerhalb einer wiederverwendbaren
technischen Ausgangsbasis.

Auf Grundlage dieser Analyse soll bewertet werden, in welchem Umfang das
Template Wiederverwendung und Standardisierung unterstützt und welches
Potenzial sich daraus für die Produktivität bei der Initialisierung und
Weiterentwicklung von Softwareprojekten ergibt. Darüber hinaus wird
untersucht, wie die vorgesehenen Projektprofile unterschiedliche
Anforderungen von Web-, Cloud- und Desktop-Anwendungen abbilden. Die Bewertung
erfolgt anhand der im Projekt vorliegenden Strukturen, Implementierungen und
Verifikationsnachweise; quantitative Effekte werden nur berücksichtigt, sofern
hierfür belastbare Daten verfügbar sind.

Als anwendungsbezogenes Ergebnis sollen geeignete Einsatzfelder und relevante
Auswahlkriterien für zukünftige Kundenprojekte abgeleitet werden. Zugleich
sind technische Grenzen, externe Voraussetzungen und projektspezifisch
verbleibende Aufgaben herauszuarbeiten. Nicht Ziel der Untersuchung ist es,
das Template als universelle Lösung für sämtliche Softwareklassen einzuordnen
oder eine produktspezifische Facharchitektur vorzugeben.
